% GENERATED FILE. Edit canonical lessons, then run npm run generate:books.
% canonical-source-hash: fnv1a64:4ee3581752bed773
% canonical-lessons: TE-C346-niscitartham, TE-C346-vivaham, TE-C346-subhalekha, TE-C346-tali, TE-C346-ahvanincu

\chapter{A Wedding}
\label{ch:te-a-wedding}
\hlchaptermodality{346}

\begin{chapteropening}
\textbf{By the end of this chapter:} \emph{I can say an engagement, a marriage, a wedding invitation, a wedding pendant and to invite.}

Complete the last lesson of chapter 346: I can say an engagement, a marriage, a wedding invitation, a wedding pendant and to invite.
\end{chapteropening}

\section[\texorpdfstring{\te{నిశ్చితార్థం} (niścitārthaṁ)}{నిశ్చితార్థం (niścitārthaṁ)}]{\te{నిశ్చితార్థం} (niścitārthaṁ) — an engagement (to marry)}

\label{lesson:TE-C346-niscitartham}



\begin{quote}
Before the new one: say the Telugu for a death anniversary, then the Telugu for a will.
\end{quote}

\subsection*{You'll want to know: \te{నిశ్చితార్థం}}
\textbf{\te{నిశ్చితార్థం}} — \emph{niścitārthaṁ} — \textquotedblleft{}an engagement (to marry)\textquotedblright{}.

\begin{quote}
\te{వాళ్ళ} \te{నిశ్చితార్థం} \te{గత} \te{నెల} \te{జరిగింది}.
\end{quote}

\begin{quote}
\emph{vāḷḷa niścitārthaṁ gata nela jarigindi.}
\end{quote}

\begin{quote}
Their engagement took place last month.
\end{quote}

\te{నిశ్చిత}, settled, and \te{అర్థం}, a matter: the day the match is settled.

\begin{cousinweb}
\textbf{\te{నిశ్చితార్థం}} (\emph{niścitārthaṁ}) is borrowed from Sanskrit: \emph{niścita}, \textquotedblleft{}ascertained, settled, decided\textquotedblright{} (Monier-Williams), and \emph{artha}, \textquotedblleft{}a matter, an aim\textquotedblright{}.
\end{cousinweb}

\begin{grammarlens}[title={What you've built}]
One more word for telling what happened.
\end{grammarlens}

\subsection*{Your turn}
\begin{itemize}
\raggedright
  \item \emph{Say it:} \emph{niścitārthaṁ}
  \item \emph{Say it:} \emph{niścitārthaṁ}, once more
  \item \emph{Say it:} \emph{vāḷḷa niścitārthaṁ gata nela jarigindi.} — Their engagement took place last month.
  \item \emph{Recall:} say the Telugu for a death anniversary, then the Telugu for a will, then say \emph{niścitārthaṁ} again
\end{itemize}

\begin{tcolorbox}[breakable,colback=teal!4,colframe=teal!35!black,title={Before you move on}]
What does \te{నిశ్చితార్థం} mean? (\textquotedblleft{}An engagement (to marry)\textquotedblright{}.) Say it once more.
\end{tcolorbox}
\section[\texorpdfstring{\te{వివాహం} (vivāhaṁ)}{వివాహం (vivāhaṁ)}]{\te{వివాహం} (vivāhaṁ) — a marriage; a wedding}

\label{lesson:TE-C346-vivaham}



\begin{quote}
Before the new one: say the Telugu for a will, then the Telugu for an engagement.
\end{quote}

\subsection*{You'll want to know: \te{వివాహం}}
\textbf{\te{వివాహం}} — \emph{vivāhaṁ} — \textquotedblleft{}a marriage; a wedding\textquotedblright{}.

\begin{quote}
\te{వాళ్ళ} \te{వివాహం} \te{గుడిలో} \te{జరిగింది}.
\end{quote}

\begin{quote}
\emph{vāḷḷa vivāhaṁ guḍilō jarigindi.}
\end{quote}

\begin{quote}
Their marriage took place in the temple.
\end{quote}

A formal word; in speech \te{పెళ్ళి} is usual.

\begin{cousinweb}
\textbf{\te{వివాహం}} (\emph{vivāhaṁ}) is borrowed from Sanskrit \emph{vivāha}, \textquotedblleft{}marriage\textquotedblright{} (CDIAL 11920).
\end{cousinweb}

\begin{grammarlens}[title={What you've built}]
One more word for telling what happened.
\end{grammarlens}

\subsection*{Your turn}
\begin{itemize}
\raggedright
  \item \emph{Say it:} \emph{vivāhaṁ}
  \item \emph{Say it:} \emph{vivāhaṁ}, once more
  \item \emph{Say it:} \emph{vāḷḷa vivāhaṁ guḍilō jarigindi.} — Their marriage took place in the temple.
  \item \emph{Recall:} say the Telugu for a will, then the Telugu for an engagement, then say \emph{vivāhaṁ} again
\end{itemize}

\begin{tcolorbox}[breakable,colback=teal!4,colframe=teal!35!black,title={Before you move on}]
What does \te{వివాహం} mean? (\textquotedblleft{}A marriage; a wedding\textquotedblright{}.) Say it once more.
\end{tcolorbox}
\section[\texorpdfstring{\te{శుభలేఖ} (śubhalēkha)}{శుభలేఖ (śubhalēkha)}]{\te{శుభలేఖ} (śubhalēkha) — a wedding invitation}

\label{lesson:TE-C346-subhalekha}



\begin{quote}
Before the new one: say the Telugu for an engagement, then the Telugu for a marriage.
\end{quote}

\subsection*{You'll want to know: \te{శుభలేఖ}}
\textbf{\te{శుభలేఖ}} — \emph{śubhalēkha} — \textquotedblleft{}a wedding invitation\textquotedblright{}.

\begin{quote}
\te{మా} \te{బంధువులు} \te{శుభలేఖ} \te{పంపించారు}.
\end{quote}

\begin{quote}
\emph{mā bandhuvulu śubhalēkha pampiñcāru.}
\end{quote}

\begin{quote}
Our relatives sent a wedding invitation.
\end{quote}

\te{శుభ}, auspicious, and \te{లేఖ}, a letter.

\begin{cousinweb}
\textbf{\te{శుభలేఖ}} (\emph{śubhalēkha}) is a compound borrowed from Sanskrit: \emph{śubha}, \textquotedblleft{}bright, auspicious\textquotedblright{} (CDIAL 12532), and \emph{lekhā}, \textquotedblleft{}a line, a stroke\textquotedblright{} (Monier-Williams).
\end{cousinweb}

\begin{grammarlens}[title={What you've built}]
One more word for telling what happened.
\end{grammarlens}

\subsection*{Your turn}
\begin{itemize}
\raggedright
  \item \emph{Say it:} \emph{śubhalēkha}
  \item \emph{Say it:} \emph{śubhalēkha}, once more
  \item \emph{Say it:} \emph{mā bandhuvulu śubhalēkha pampiñcāru.} — Our relatives sent a wedding invitation.
  \item \emph{Recall:} say the Telugu for an engagement, then the Telugu for a marriage, then say \emph{śubhalēkha} again
\end{itemize}

\begin{tcolorbox}[breakable,colback=teal!4,colframe=teal!35!black,title={Before you move on}]
What does \te{శుభలేఖ} mean? (\textquotedblleft{}A wedding invitation\textquotedblright{}.) Say it once more.
\end{tcolorbox}
\section[\texorpdfstring{\te{తాళి} (tāḷi)}{తాళి (tāḷi)}]{\te{తాళి} (tāḷi) — a wedding pendant}

\label{lesson:TE-C346-tali}



\begin{quote}
Before the new one: say the Telugu for a marriage, then the Telugu for a wedding invitation.
\end{quote}

\subsection*{You'll want to know: \te{తాళి}}
\textbf{\te{తాళి}} — \emph{tāḷi} — \textquotedblleft{}a wedding pendant\textquotedblright{}.

\begin{quote}
\te{పెళ్ళికొడుకు} \te{పెళ్ళికూతురు} \te{మెడలో} \te{తాళి} \te{కట్టాడు}.
\end{quote}

\begin{quote}
\emph{peḷḷikoḍuku peḷḷikūturu meḍalō tāḷi kaṭṭāḍu.}
\end{quote}

\begin{quote}
The bridegroom tied the wedding pendant round the bride's neck.
\end{quote}

\te{తాళి} \te{కట్టు}, to tie the tali: the moment a Telugu wedding is sealed.

\begin{cousinweb}
\textbf{\te{తాళి}} (\emph{tāḷi}) is native Dravidian: DEDR 3175 lists Telugu \emph{tāli}, the small piece of gold the bridegroom ties at the neck of the bride, with Tamil \emph{tāli}.
\end{cousinweb}

\begin{grammarlens}[title={What you've built}]
One more word for telling what happened.
\end{grammarlens}

\subsection*{Your turn}
\begin{itemize}
\raggedright
  \item \emph{Say it:} \emph{tāḷi}
  \item \emph{Say it:} \emph{tāḷi}, once more
  \item \emph{Say it:} \emph{peḷḷikoḍuku peḷḷikūturu meḍalō tāḷi kaṭṭāḍu.} — The bridegroom tied the wedding pendant round the bride's neck.
  \item \emph{Recall:} say the Telugu for a marriage, then the Telugu for a wedding invitation, then say \emph{tāḷi} again
\end{itemize}

\begin{tcolorbox}[breakable,colback=teal!4,colframe=teal!35!black,title={Before you move on}]
What does \te{తాళి} mean? (\textquotedblleft{}A wedding pendant\textquotedblright{}.) Say it once more.
\end{tcolorbox}
\section[\texorpdfstring{\te{ఆహ్వానించు} (āhvāniñcu)}{ఆహ్వానించు (āhvāniñcu)}]{\te{ఆహ్వానించు} (āhvāniñcu) — to invite}

\label{lesson:TE-C346-ahvanincu}



\begin{quote}
Before the new one: say the Telugu for a wedding invitation, then the Telugu for a wedding pendant.
\end{quote}

\subsection*{You'll want to know: \te{ఆహ్వానించు}}
\textbf{\te{ఆహ్వానించు}} — \emph{āhvāniñcu} — \textquotedblleft{}to invite\textquotedblright{}.

\begin{quote}
\te{వాళ్ళు} \te{మమ్మల్ని} \te{పెళ్ళికి} \te{ఆహ్వానించారు}.
\end{quote}

\begin{quote}
\emph{vāḷḷu mammalni peḷḷiki āhvāniñcāru.}
\end{quote}

\begin{quote}
They invited us to the wedding.
\end{quote}

Past (I …): \te{ఆహ్వానించాను}, \emph{āhvāniñcānu}.

You know the noun \te{ఆహ్వానం}, an invitation; the event takes -\te{కి}.

\begin{cousinweb}
\textbf{\te{ఆహ్వానించు}} (\emph{āhvāniñcu}) is built on Sanskrit \emph{āhvāna}, \textquotedblleft{}calling, an invitation\textquotedblright{} (Monier-Williams), with the Telugu verb ending \emph{-iñcu}.
\end{cousinweb}

\begin{grammarlens}[title={What you've built}]
One more word for telling what happened.
\end{grammarlens}

\subsection*{Your turn}
\begin{itemize}
\raggedright
  \item \emph{Say it:} \emph{āhvāniñcu}
  \item \emph{Say it:} \emph{āhvāniñcu}, once more
  \item \emph{Say it:} \emph{vāḷḷu mammalni peḷḷiki āhvāniñcāru.} — They invited us to the wedding.
  \item \emph{Recall:} say the Telugu for a wedding invitation, then the Telugu for a wedding pendant, then say \emph{āhvāniñcu} again
\end{itemize}

\begin{tcolorbox}[breakable,colback=teal!4,colframe=teal!35!black,title={Before you move on}]
What does \te{ఆహ్వానించు} mean? (\textquotedblleft{}To invite\textquotedblright{}.) Say it once more.
\end{tcolorbox}
